% ======================================================================
%  Section 4.4.2 --- Local Hermite structure measure.  Second correction.
%
%  A. A SENTENCE HAS BEEN BROKEN IN TWO BY THE EQUATION. As it stands the
%     text reads
%         "... the first derivatives of the Gaussian of width sigma.
%          [equation]
%          Their responses, are accordingly a Gaussian-smoothed image
%          gradient ..."
%     The equation was moved out of the sentence that introduced it,
%     leaving a stray comma and a fragment. This is visible in the output.
%     Repaired below.
%
%  B. "must do so from the Hermite bank the scheme ALREADY COMPUTES" and
%     "is obtained from kernels the scheme ALREADY EVALUATES" -- the fourth
%     and fifth occurrences of the claim corrected in 4.1, 4.3.3 and 4.4.1.
%     Both are mine from the earlier draft. The first-order kernels are not
%     among those the feature vector uses.
%
%  C. NOTATION COLLISION, and it is in this subsection that the two meet.
%     D_{n,m}^{sigma} is a Gauss-Hermite KERNEL here, while D, D^T and D^H
%     are the WEIGHT MATRICES of 4.2.2, 4.3.2 and 4.4.3. Same bold letter,
%     two unrelated objects, both in Section 4.4. Since the kernel notation
%     comes from Chapter 3 and is fixed by the implementation, the weight
%     matrix is the one to rename:
%         D  ->  W        W = diag(w_1,...,w_P)
%         D^T -> W^T      D^H -> W^H
%     It touches 4.2.2, 4.3.2, 4.4.3 and 4.6.1 and nothing else, and it
%     also removes the transpose ambiguity flagged earlier.
%
%  D. LABEL CONFLICT. This subsection refers to Chapter 3's locality result
%     as sec:locality, twice. Sections 4.1, 4.3.3 and 4.4.1 refer to the
%     same result as sec:hermite-properties. One of the two will print as
%     Section ??. Pick whichever exists in Chapter 3 and use it everywhere.
%
%  E. A %%% COMMENT IS EMBEDDED INSIDE A SENTENCE, between
%     Section~\ref{sec:ablation-measure} and "records what happens".
%     It compiles, because % eats the newline, but it should not be there.
% ======================================================================

\subsection{Local Hermite structure measure}
\label{sec:hermite-gate}

%%% (B) first occurrence corrected: "already computes" -> the weaker and
%%% true statement that the measure is drawn from the same bank.
The measure must report how much exploitable structure the image contains at
a given position, and must do so within the Hermite representation the scheme
already uses. Not every response of that bank is suitable. Writing
$\mathbf{D}_{n,m}^{\sigma}$ for the separable Gauss--Hermite kernel of orders
$(n,m)$ at scale $\sigma$, as defined in Section~\ref{sec:hermite-bank}, the
feature vector of Chapter~\ref{ch:herpvs} uses the sum of these kernels over
all orders $n,m\le4$, and that sum is low-pass dominated: its gain at zero
spatial frequency equals its $L_{1}$ norm to three decimal places. Its
response is therefore close to a local average of the image, and would report
local \emph{intensity} rather than local structure.
Section~\ref{sec:ablation-measure} records what happens when it is used.
%%% (E) the stray %%% comment that sat inside this sentence is removed.

%%% (A) The sentence is made whole again: subject, equation, predicate, in
%%% that order, with the equation as the object of "Their responses".
The first-order kernels do not have this defect. Since the Hermite function
of order one is odd, $\mathbf{D}_{1,0}^{\sigma}$ and
$\mathbf{D}_{0,1}^{\sigma}$ integrate to zero over their support exactly, and
are, up to a scale factor, the first derivatives of the Gaussian of width
$\sigma$. Their responses,
\begin{equation}
    g_{x}^{(k)}(p)=\bigl(\mathbf{D}_{1,0}^{\sigma_{k}}\ast\mathbf{I}\bigr)(p),
    \qquad
    g_{y}^{(k)}(p)=\bigl(\mathbf{D}_{0,1}^{\sigma_{k}}\ast\mathbf{I}\bigr)(p),
    \label{eq:first-order-responses}
\end{equation}
are accordingly a Gaussian-smoothed image gradient at scale $\sigma_{k}$, and
the local Hermite structure measure is defined as their combined magnitude
over the three coarser scales of the bank,
\begin{equation}
    G(p)=\sqrt{\sum_{k=2}^{4}\Bigl[\bigl(g_{x}^{(k)}(p)\bigr)^{2}
                                  +\bigl(g_{y}^{(k)}(p)\bigr)^{2}\Bigr]}.
    \label{eq:structG}
\end{equation}

%%% (B) second occurrence corrected: the trailing clause "and is obtained
%%% from kernels the scheme already evaluates" is deleted. The sentence is
%%% stronger without it, and 4.4.1 now states the cost honestly.
Being built from zero-mean kernels, $G$ is invariant to any additive change
in illumination and vanishes on a region of uniform intensity whatever that
intensity is --- the property on which Section~\ref{sec:occlusion-mechanism}
depends. It is the band-pass character of
Section~\ref{sec:hermite-properties} that makes this possible; what
\eqref{eq:structG} contributes beyond a generic image-gradient measure is not
the form of the operator, which is a classical one, but that it is evaluated
at precisely the analysis scales $\sigma_{2},\sigma_{3},\sigma_{4}$ that
define the features the controller acts on, and that its magnitude is, by the
argument of Section~\ref{sec:hermite-principle}, the image-dependent factor
of the norm of the corresponding row of $\mathbf{L}_{R}$.
%%% The replacement clause ties this subsection to the interaction-matrix
%%% argument added to 4.4.1, which is what actually justifies the choice.

\paragraph{The two measures compared.}
Figure~\ref{fig:structure-maps} shows the two forms of~\eqref{eq:structG} for
the target of Figure~\ref{fig:setup}(a). The difference is not subtle. Built
from the summed kernels, $G$ reproduces the photograph: the mean magnitude of
its gradient is $6.2\%$ of its mean level and $62.8\%$ of positions lie
within a tenth of their own $9\times9$ local mean, which is what a smoothed
copy of the intensity looks like. Built from the first-order kernels, the same
quantities become $26.2\%$ and $18.8\%$, and the map traces the contours of
the scene while falling to zero inside regions of uniform intensity, whether
those regions are light or dark.

%%% Sharpened. "Only the odd orders taken alone" is loose: D_{1,1} is even
%%% in neither variable and also integrates to zero, while D_{2,0} does not.
%%% The governing condition is that at least one order be odd.
It is worth noting that excluding the $(0,0)$ term from the summed kernels is
not sufficient. Measured, that reduces the ratio of DC gain to $L_{1}$ norm
only from $1.000$ to $0.810$, since a kernel $\mathbf{D}_{n,m}^{\sigma}$
integrates to zero if and only if at least one of $n$ and $m$ is odd, and the
terms with both orders even continue to carry a DC component. Only kernels
satisfying that condition give a measure exactly free of it.

\paragraph{Exclusion of the finest scale.}
The scale $\sigma_{1}$ is omitted from~\eqref{eq:structG}. The reason is one
of discretisation rather than of scale selection: at $\sigma_{1}=0.5$ the
support of the kernel is comparable to the sampling interval, so on the
$9\times9$ grid the discrete first-order response reduces to a difference over
two or three taps and is dominated by pixel noise rather than by image
structure. The noise suppression established in
Section~\ref{sec:hermite-properties}, which follows from the Gaussian envelope
being wide relative to the sampling interval, does not yet apply at that
scale. The three coarser scales are retained, and no attempt is made here to
optimise the set; the unequal contribution of the four scales to the feature
error, quantified in Section~\ref{sec:scale-energy}, indicates that such an
optimisation is warranted, though it is not attempted here.

\begin{figure}[H]
    \centering
    \begin{minipage}[b]{0.47\textwidth}
        \centering
        \includegraphics[width=\textwidth]{Figs/D/gmap_v3_600_nominal_intensity.png}

        \vspace{2mm}
        \textbf{(a)} Summed kernels: $G$ tracks intensity.
    \end{minipage}
    \hfill
    \begin{minipage}[b]{0.47\textwidth}
        \centering
        \includegraphics[width=\textwidth]{Figs/D/gmap_v3_600nominal.png}

        \vspace{2mm}
        \textbf{(b)} First-order kernels: $G$ tracks structure.
    \end{minipage}

    \caption{The local structure measure~\eqref{eq:structG} displayed at its
    image positions, with brightness proportional to $G(p)$ and scaled so
    that three times the image mean saturates. (a) Built from the summed
    kernels $\mathbf{D}_{n,m}$, $n,m\leq4$, whose DC gain equals their
    $L_{1}$ norm: the measure reproduces the target image. (b) Built from
    the first-order kernels of~\eqref{eq:first-order-responses}, which
    integrate to zero: the measure traces the contours of the scene and
    vanishes in regions of uniform intensity. Both panels are taken at
    iteration~600 of the nominal run and differ by one line of the
    implementation.}
    \label{fig:structure-maps}
\end{figure}

% ======================================================================
%  F. SMALLER POINTS
%
%  - \paragraph{The two measures compared:} had a colon inside the braces.
%    LaTeX's \paragraph supplies its own terminal punctuation; a full stop
%    is the convention. Corrected in both paragraph headings.
%
%  - \begin{figure}[H] requires \usepackage{float}. If the preamble does
%    not load it the document will not compile. [htbp] is the safer choice
%    unless you specifically need the figure pinned here.
%
%  - The caption said the panels are "from the same iteration of the same
%    run configuration". They are from the same iteration of two DIFFERENT
%    configurations -- that is the whole point of the comparison. Reworded
%    to say iteration 600 of the nominal run, differing by one line.
%
%  - Figure file names are inconsistent: gmap_v3_600_nominal_intensity.png
%    against gmap_v3_600nominal.png. Harmless, but rename one so the pair
%    sorts together in Figs/D.
%
%  LABELS THIS SUBSECTION DEPENDS ON -- confirm each exists:
%    sec:hermite-bank         Chapter 3, where D_{n,m}^sigma is defined
%    sec:hermite-properties   Chapter 3, locality and band-pass (see D)
%    sec:hermite-principle    4.4.1, added in the previous correction
%    sec:occlusion-mechanism  4.4.3, the paragraph on the occluded region
%    sec:ablation-measure     4.4.4
%    sec:scale-energy         the per-scale energy table
%    sec:future               future work
%    fig:setup                Figure 4.2
% ======================================================================
